\subsection{Flatness of quotient modules}

PROPOSITION 4. Let E be a right A-module. The three following properties are equivalent:

\begin{enumerate}
\def\labelenumi{(\alph{enumi})}
\item
  E is flat.
\item
  For every exact sequence of right A-modules of the form
\end{enumerate}

\[
0 \longrightarrow \mathrm{G} \stackrel {v} {\longrightarrow} \mathrm{H} \stackrel {w} {\longrightarrow} \mathrm{E} \longrightarrow 0\tag{1}
\]

and every left A-module F, the sequence

\[
0 \longrightarrow \mathrm{G} \otimes_ {\mathrm{A}} \mathrm{F} \xrightarrow {v \otimes 1} \mathrm{H} \otimes_ {\mathrm{A}} \mathrm{F} \xrightarrow {w \otimes 1} \mathrm{E} \otimes_ {\mathrm{A}} \mathrm{F} \longrightarrow 0\tag{2}
\]

is exact.

\begin{enumerate}
\def\labelenumi{(\alph{enumi})}
\setcounter{enumi}{2}
\tightlist
\item
  There exists an exact sequence (1), where H is flat, such that the sequence (2) is exact for every left A-module F of the form \(A_{s}/a\) , where a is a finitely generated left ideal of A.
\end{enumerate}

We show first that (a) implies (b). The left A-module F is isomorphic to a quotient of a free module (Algebra, Chapter II, § 1, no. 11, Proposition 20); in other words, we have an exact sequence

\[
0 \longrightarrow \mathrm{R} \stackrel {i} {\longrightarrow} \mathrm{L} \stackrel {p} {\longrightarrow} \mathrm{F} \longrightarrow 0
\]

where L is free. Consider the diagram

\[
\begin{array}{c} \mathrm{G} \otimes \mathrm{R} \xrightarrow {v \otimes 1 _ {\mathrm{R}}} \mathrm{H} \otimes \mathrm{R} \xrightarrow {w \otimes 1 _ {\mathrm{R}}} \mathrm{E} \otimes \mathrm{R} \\ \downarrow_ {1 _ {\mathrm{G}} \otimes i} \\ \mathrm{G} \otimes \mathrm{L} \xrightarrow {v \otimes 1 _ {\mathrm{L}}} \mathrm{H} \otimes \mathrm{L} \xrightarrow {w \otimes 1 _ {\mathrm{L}}} \mathrm{E} \otimes \mathrm{L} \\ \downarrow_ {1 _ {\mathrm{G}} \otimes p} \\ \mathrm{G} \otimes \mathrm{F} \xrightarrow {v \otimes 1 _ {\mathrm{F}}} \mathrm{H} \otimes \mathrm{F} \end{array}\tag{3}
\]

It follows immediately that this diagram is commutative, and its rows and columns are exact by Lemma 1 of no. 1; moreover, as \(\mathbf{1}_{\mathbf{G}} \otimes p\) and \(\mathbf{1}_{\mathbf{H}} \otimes p\) are surjective (no. 1, Lemma 1), we have \(\mathbf{G} \otimes \mathbf{F} = \operatorname{Coker}(\mathbf{1}_{\mathbf{G}} \otimes i)\) ,

\[
\mathbf {H} \otimes \mathbf {F} = \operatorname{Coker} (\mathbf {1} _ {\mathbf {H}} \otimes i);
\]

\(w \otimes l_{R}\) is surjective (no. 1, Lemma 1); finally, as L is free and hence flat, \(v \otimes l_{L}\) is injective. Thus the snake diagram (§ 1, no. 4, Proposition 2, (iii)) can be applied to prove the existence of an exact sequence

\[
\operatorname{Ker} \left(\mathrm{l} _ {\mathrm{H}} \otimes i\right) \longrightarrow \operatorname{Ker} \left(\mathrm{l} _ {\mathrm{E}} \otimes i\right) \xrightarrow {d} \mathrm{G} \otimes \mathrm{F} \xrightarrow {v \otimes \mathrm{l} _ {\mathrm{F}}} \mathrm{H} \otimes \mathrm{F}.\tag{4}
\]

This being so, if E is flat, \(1_{E} \otimes i\) is injective, in other words \(\operatorname{Ker}(1_{\mathbf{E}} \otimes i) = 0\) , and the exact sequence (4) shows that \(v \otimes 1_{F}\) is injective, hence the sequence (2) is exact (taking account of Lemma 1 of no. 1).

As (b) obviously implies (c), it remains to prove that (c) implies (a). Consider the diagram (3) in the case \(\mathbf{R} = \mathfrak{a}\) , \(\mathbf{L} = \mathbf{A}_s\) , \(\mathbf{F} = \mathbf{A}_s / \mathfrak{a}\) and apply the exact sequence (4). By hypothesis \(v \otimes 1_{\mathbb{F}}\) is injective, hence \(\operatorname{Im}(d) = 0\) ; moreover, as H is flat, we have \(\operatorname{Ker}(1_{\mathbb{H}} \otimes i) = 0\) ; the exactness of the sequence (4) then implies \(\operatorname{Ker}(1_{\mathbb{E}} \otimes i) = 0\) , in other words, \(1_{\mathbb{E}} \otimes i\) is injective and this proves that E is flat (no. 3, Remark 1).

PROPOSITION 5. Let \(0 \to E' \xrightarrow{v} E \xrightarrow{w} E'' \to 0\) be an exact sequence of right A-modules. Suppose \(E''\) is flat. Then, for \(E\) to be flat it is necessary and sufficient that \(E'\) be flat.

Let \(u: F' \to F\) be an injective homomorphism of left A-modules. Consider the diagram

\[
\begin{array}{c} \mathrm{E} ^ {\prime} \otimes \mathrm{F} ^ {\prime} \xrightarrow {v \otimes 1 _ {\mathrm{F} ^ {\prime}}} \mathrm{E} \otimes \mathrm{F} ^ {\prime} \xrightarrow {w \otimes 1 _ {\mathrm{F} ^ {\prime}}} \mathrm{E} ^ {\prime \prime} \otimes \mathrm{F} ^ {\prime} \\ 1 _ {\mathrm{E} ^ {\prime}} \otimes u \Bigg \downarrow \qquad \qquad \qquad \qquad \qquad \qquad \qquad \qquad \qquad \qquad \qquad \qquad \qquad \qquad \qquad \qquad \qquad \qquad \qquad \qquad \qquad \qquad \qquad \qquad \qquad \qquad \qquad \qquad \qquad \qquad \qquad \qquad \qquad \qquad \\ \mathrm{E} ^ {\prime} \otimes \mathrm{F} \xrightarrow [ v \otimes 1 _ {\mathrm{F}} ]{} \mathrm{E} \otimes \mathrm{F} \xrightarrow [ w \otimes 1 _ {\mathrm{F}} ]{} \mathrm{E} ^ {\prime \prime} \otimes \mathrm{F} \end{array}
\]

It is commutative and its rows are exact (no. 1, Lemma 1). Since \(\mathbf{E}''\) is flat,

\(1_{\mathbf{E}^{\prime \prime}} \otimes u\) is injective; moreover, Proposition 4 proves that \(v \otimes 1_{\mathbf{F}^{\prime}}\) and \(v \otimes 1_{\mathbf{F}}\) are injective. This being so, if E is flat, \(1_{\mathbf{E}} \otimes u\) is injective, hence also

\[
\left(1 _ {\mathbf {E}} \otimes u\right) \circ (v \otimes 1 _ {\mathbf {F} ^ {\prime}}) = (v \otimes 1 _ {\mathbf {F}}) \circ \left(1 _ {\mathbf {E} ^ {\prime}} \otimes u\right);
\]

it follows that \(\mathbf{1}_{\mathbf{E}'} \otimes u\) is injective and consequently \(\mathbf{E}'\) is flat. Conversely, if \(\mathbf{E}'\) is flat, \(\mathbf{1}_{\mathbf{E}'} \otimes u\) is injective; then it follows from §1, no. 4, Corollary 1 to Proposition 2 that \(\mathbf{1}_{\mathbf{E}} \otimes u\) is injective and so \(\mathbf{E}\) is flat.

\textbf{Remarks}\par

\begin{enumerate}
\def\labelenumi{(\arabic{enumi})}
\item
  It is possible for E and \(E'\) to be flat without \(E''\) being so, as the example of the Z-modules \(E = Z\) , \(E' = nZ\) , \(E'' = Z/nZ\) ( \(n \geqslant 2\) ), shows.
\item
  A submodule of a flat module is not necessarily a flat module (Exercise 3).
\end{enumerate}

